\documentclass[11pt]{article}

\usepackage[margin=1in]{geometry}
\usepackage{amsmath,amssymb,amsthm,mathtools}
\usepackage{enumitem}
\usepackage{booktabs}
\usepackage{tabularx}
\usepackage{microtype}
\usepackage[hidelinks]{hyperref}
\usepackage{xurl}

\newtheorem{definition}{Definition}
\newtheorem{lemma}{Lemma}
\newtheorem{remark}{Remark}

\title{Release Obligations and Closure Ledgers for Successor Maintenance:\\[0.5ex]
\large Required Roles, Row Satisfaction, and Completion Tokens}
\author{}
\date{Unified archive note v0.05 (2026-03-21; archive v489)}

\begin{document}
\maketitle

\begin{abstract}
A successor release often needs one public answer to two different questions: which outputs are required after the observed drift, and whether the shipped package actually satisfies those duties.
Those questions should not collapse into one opaque token.
This note separates them cleanly.
It defines an \emph{obligation profile} that names the required publication roles for each winning action class, a \emph{closure ledger} that checks those roles row-by-row against concrete shipped artifacts, and a \emph{closure verdict} that compresses the ledger to one completion token only after the role rows have been bound.
The result is a small auditable bridge between release classification and downstream public wording.
\end{abstract}

\paragraph{Non-redundancy note.}
Synthesis~18 owns the drift taxonomy and action classes.\cite{series:planstateequiv}
Synthesis~20 owns replay mechanics and replay verdicts.\cite{series:auditorreplay}
Synthesis~32 owns the stable release spine and stage order.\cite{series:releasespine}
Synthesis~31 owns the downstream sentence-reuse suffix once the closure token has already been established.\cite{series:downstreamnormalforms}
This note owns only the obligation$\to$closure bridge.

\paragraph{Scope.}
This is not a new receipt primitive and not a new certification mechanism.
It is a publication-interface note for successor maintenance after drift classification has already occurred.

\paragraph{Imports.}
Contract-object vocabulary is imported from Synthesis~0.\cite{series:contractobjects}
Release-drift classes and obligation triggers are imported from Synthesis~18.\cite{series:planstateequiv}
The stable stage order is imported from Synthesis~32.\cite{series:releasespine}
The answer-side terminal reuse / cutover discipline is imported from Synthesis~54 and Synthesis~74--75, and the concrete worked instantiation is imported from Synthesis~17.\cite{series:challengeanswerreviewmenus,series:pairedterminalcitationdiscipline,series:pairedterminalcutoverrules,series:workedexample}

\paragraph{Exports.}
This note exports (i) a required-role row for successor publication duties, (ii) a closure-ledger row that binds that role to concrete shipped artifacts, (iii) a completeness rule for turning the ledger into one closure token, (iv) a minimal citation boundary separating policy, row-level checking, and compressed public status, and (v) a local-reopen rule saying that later terse papers should reopen this note only when the exact obligation profile, closure-ledger row, or closure verdict is itself live. Otherwise, once the checked-answer handoff has already been fixed, later terse papers should default to the answer-side terminal citation normal form in Synthesis~54 together with the paired terminal discipline and paired cutover rule in Synthesis~74--75.\cite{series:challengeanswerreviewmenus,series:pairedterminalcitationdiscipline,series:pairedterminalcutoverrules}

\paragraph{Later-terse-paper citation posture.}
When a later short paper only needs to answer \emph{which release-obligation / closure packet was declared for one maintained successor package?}, it should stop at the minimal public obligation-and-closure card below plus the tiny row-satisfaction drill. It should widen only when the live issue becomes drift classification (Synthesis~18), replay mechanics (Synthesis~20), summary-stage status compression or wrapper prose (Synthesis~34), release-stage ownership (Synthesis~32), or a checked-answer terminal handoff (Synthesis~54 together with Synthesis~74--75).\cite{series:planstateequiv,series:auditorreplay,series:statusenvelopes,series:releasespine,series:challengeanswerreviewmenus,series:pairedterminalcitationdiscipline,series:pairedterminalcutoverrules}

\section{Three layers, not one}

\begin{definition}[Required-role row]
For one action class $a$, a \emph{required-role row}
\[
R=(r,b,\Gamma)
\]
consists of a publication role label $r$, a required binding description $b$, and a finite candidate set $\Gamma$ of concrete artifact paths that may satisfy that role in a shipped package.
The row says what must be present if $a$ wins.
It does not say whether the current package is complete.
\end{definition}

\begin{definition}[Closure-ledger row]
For one required-role row $R=(r,b,\Gamma)$ and one concrete shipped package $P$, a \emph{closure-ledger row}
\[
L_R(P)=(r,b,\Gamma,B_R(P),s_R(P))
\]
records the subset $B_R(P)\subseteq \Gamma$ actually bound in the package and one row status
\[
s_R(P)\in\{\textsf{satisfied},\textsf{missing}\}.
\]
The row is satisfied exactly when $B_R(P)\neq \varnothing$.
\end{definition}

\begin{definition}[Closure ledger]
For one winning action class $a$ and package $P$, the \emph{closure ledger}
\[
\Lambda_a(P)=\{L_R(P): R \in \mathcal R(a)\}
\]
contains one closure-ledger row for every required-role row in the obligation profile of $a$.
The ledger owns row-level package checking.
It is the smallest object that can answer ``which required role is missing?'' without reopening drift classification or replay.
\end{definition}

\begin{definition}[Closure verdict]
A \emph{closure verdict} for $a$ and $P$ is the pair
\[
V_a(P)=(\Lambda_a(P), c_a(P))
\]
where
\[
c_a(P)=\textsf{complete}
\quad\text{iff}\quad
s_R(P)=\textsf{satisfied}\ \text{for every row in }\Lambda_a(P).
\]
The verdict may publish the token $c_a(P)$ and the bound outputs that witness it, but it does not replace the ledger.
\end{definition}

\begin{remark}[Why the separation matters]
An obligation profile answers a policy question: \emph{what would be required if this class wins?}
A closure ledger answers an audit question: \emph{which required roles are actually satisfied by this package?}
A closure verdict answers a public-status question: \emph{is the package complete?}
Collapsing those three layers into one object makes later papers either over-cite policy or under-cite row-level evidence.
\end{remark}

\section{Completeness and citation boundaries}

\begin{lemma}[Ledger completeness rule]
Let $a$ be the winning action class for one base$\to$successor comparison and let $P$ be the shipped successor package.
Then the closure token $c_a(P)$ is auditable from the closure ledger alone:
\[
c_a(P)=\textsf{complete}
\quad\Longleftrightarrow\quad
\forall R\in\mathcal R(a),\ B_R(P)\neq\varnothing.
\]
\end{lemma}

\begin{proof}
By definition, the closure ledger contains one row per required role of the winning action class.
Each row status is \textsf{satisfied} exactly when at least one allowed artifact candidate is actually bound in the package.
The closure token is \textsf{complete} exactly when every row is satisfied.
So the token is just a compression of the row statuses already present in the ledger.
\end{proof}

\begin{remark}[Minimal citation rule]
Use the obligation profile when the claim is about \emph{what should be published} for a winning class.
Use the closure ledger when the claim is about \emph{which role rows are satisfied or missing} in one concrete package.
Use the closure verdict only when the claim is merely the compressed completion token itself.
This is the smallest rule that stops package-completeness prose from erasing row-level evidence. Later terse papers should otherwise default to the answer-side terminal citation normal form and paired cutover rule once the checked-answer route has already imported this closure basis.\cite{series:challengeanswerreviewmenus,series:pairedterminalcitationdiscipline,series:pairedterminalcutoverrules}
\end{remark}

\paragraph{Later-terse-paper citation posture.}
Later terse papers should cite Synthesis~33 only when the live issue is the required-role row set for one winning action class, the concrete closure-ledger binding or satisfaction status of one role row, or the final completion verdict for one shipped package. They should stop at Synthesis~18 when the live issue is the imported drift taxonomy or winning action class, stop at Synthesis~20 when the live issue is replay mechanics or replay-derived verdict fields, stop at Synthesis~32 when the question has already widened to release-stage order or stage-local handoff ownership, and stop at Synthesis~34 when the closure basis has already been compressed to the narrow public-status token or outward-facing lineage wrapper seen by readers. If the release question has already reduced to downstream summary / support / citation / quote / clause reuse, stop at Synthesis~31. Once a checked-answer handoff already imports the closure basis, later terse papers should widen no further through this note and instead default to Synthesis~54 together with Synthesis~74--75.\cite{series:planstateequiv,series:auditorreplay,series:releasespine,series:statusenvelopes,series:downstreamnormalforms,series:challengeanswerreviewmenus,series:pairedterminalcitationdiscipline,series:pairedterminalcutoverrules}


\paragraph{A minimal public obligation-and-closure card.}
\begin{table}[t]
\centering
\small
\begin{tabularx}{0.97\linewidth}{@{}p{0.31\linewidth}X@{}}
\toprule
\textbf{Public row} & \textbf{Pinned value}\\
\midrule
Winning release class & Worked Case-B successor with \texttt{classification = replay\_changed\_surfaces} from the declared release-action matrix \\
Required-role floor & Public output roles are \texttt{fresh\_release\_binding}, \texttt{fresh\_compare\_report}, \texttt{fresh\_replay\_verdict}, \texttt{fresh\_successor\_continuity\_verdict}, and \texttt{fresh\_successor\_lineage\_notice} \\
Row-level ledger basis & \texttt{example\_release\_closure\_ledger.json} binds each required role to at least one shipped artifact path and records \texttt{status = satisfied} row by row \\
Missing-role summary & The same ledger exports \texttt{missing\_publication\_roles = []}, so no required row is still open in the shipped package \\
Compressed public verdict & \texttt{example\_publication\_closure\_verdict.json} compresses that ledger to \texttt{closure\_status = complete} without replacing the row bindings \\
Escalation pointer & Widen to Synthesis~18 for the winning drift/action class, Synthesis~20 for replay mechanics, Synthesis~34 for public-status compression, and Synthesis~32 for stage-order questions \\
\bottomrule
\end{tabularx}
\caption{A minimal public obligation-and-closure card. This is the smallest public answer later terse papers can quote directly when the live issue is the declared successor-publication duty set plus its row-level completion verdict rather than replay mechanics, public-status wrapper prose, or release-stage routing.}
\label{tab:minimal-obligation-closure-card}
\end{table}

This card is intentionally non-decorative: it fixes one winning action class, one finite required-role floor, one row-by-row closure ledger, and one compressed completion token without silently re-owning drift classification, replay execution, or the public-status sentence shown to readers.\cite{series:workedexample}

\paragraph{One tiny row-satisfaction drill.}
In the worked Case-B successor, the obligation profile requires five public roles and the closure ledger binds all five of them to shipped artifacts: the release receipt, compare report, verifier report, successor continuity verdict, and successor lineage notice. Because every required row has a nonempty \texttt{bound\_artifacts} list, the ledger's \texttt{missing\_publication\_roles} field is empty and the compressed verdict may honestly publish \texttt{closure\_status = complete}. If even one required row were left without a bound artifact --- for example the successor lineage notice row --- then the honest public answer here would stay at the ledger and report that missing role rather than claiming package completeness.\cite{series:workedexample}

\begin{table}[t]
\centering
\small
\begin{tabularx}{\linewidth}{@{}l l X@{}}
\toprule
Layer & Canonical object & Question answered \\
\midrule
Policy & obligation profile & Which publication roles are required for the winning action class? \\
Row check & closure ledger & Which concrete artifact path satisfies each required role, and which rows are missing? \\
Compression & closure verdict & Is the successor package complete once those role rows are checked? \\
\bottomrule
\end{tabularx}
\caption{Policy, row-level checking, and completion tokens should stay separable.}
\label{tab:layers}
\end{table}

\section{Worked example pointer}
The maintained worked bundle ships all three layers.\cite{series:workedexample}
\path{example_release_obligation_profile.json} names the required publication roles for the winning class.
\path{example_release_closure_ledger.json} binds those role rows to the concrete shipped artifacts for the canned Case-B successor, including the allowed candidates, the artifact(s) actually bound, and a row status.
\path{example_publication_closure_verdict.json} then compresses that ledger to the completion token \texttt{complete} together with the bound output list.
That worked cut gives later papers one concrete auditable bridge to cite without forcing Synthesis~32 to own role rows or Synthesis~31 to explain package closure. In the maintained worked series spine, the same closure basis is imported directly onto each shipped audience/question route, so later terse notes can recover the governing obligation profile, closure ledger, and closure verdict without side-jumping out of the route object. Once the checked-answer handoff has already been fixed, they should then default to Synthesis~54 and Synthesis~74--75 rather than restating this obligation/closure bridge locally.\cite{series:challengeanswerreviewmenus,series:pairedterminalcitationdiscipline,series:pairedterminalcutoverrules}

\paragraph{Why this helps the series.}
The archive aims for short papers with stable reuse rules.
Owning obligation/closure semantics here lets the release-spine note stay about stage order, lets the downstream-normal-form note stay about sentence reuse, and lets the worked example stay a concrete instantiation rather than a theory surrogate.

\begin{thebibliography}{99}\small

\bibitem{series:contractobjects}
Anonymous.
\newblock Contract Objects for Anonymous-DHT Claims: Commitments, Menus, Receipts, and Release Bindings.
\newblock Series synthesis note (Synthesis~0), 2026. (This archive.)

\bibitem{series:planstateequiv}
Anonymous.
\newblock Plan and State Equivalence: Digest-Bound Replay Hooks and Drift Taxonomy.
\newblock Series synthesis note (Synthesis~18), 2026. (This archive.)

\bibitem{series:auditorreplay}
Anonymous.
\newblock Auditor Replay Recipe for Anonymous-DHT Receipts: A Minimal Checklist and Verdict Tuple.
\newblock Series synthesis note (Synthesis~20), 2026. (This archive.)

\bibitem{series:releasespine}
Anonymous.
\newblock Stable Release Spine Contracts for Successor Maintenance: Stage Ownership, Handoff Rules, and Auditable Walk-Throughs.
\newblock Series synthesis note (Synthesis~32), 2026. (This archive.)

\bibitem{series:challengeanswerreviewmenus}
Anonymous.
\newblock Challenge Answer Review Menus for Successor Maintenance: Smallest-Sufficient Referee Handoffs from Dockets, Profiles, Packets, and Ladders.
\newblock Series synthesis note (Synthesis~54), 2026. (This archive.)

\bibitem{series:pairedterminalcitationdiscipline}
Anonymous.
\newblock Paired Terminal Citation Discipline for Anonymous-DHT Receipts: Stable Answer Stops and Adjacent Request Reopenings.
\newblock Series synthesis note (Synthesis~74), 2026. (This archive.)

\bibitem{series:pairedterminalcutoverrules}
Anonymous.
\newblock Paired Terminal Cutover Rules for Anonymous-DHT Receipts: Stop, Widen, and Reopen as a Stable Terse-Paper Discipline.
\newblock Series synthesis note (Synthesis~75), 2026. (This archive.)

\bibitem{series:statusenvelopes}
Anonymous.
\newblock Successor Status Envelopes and Lineage Notices for Anonymous-DHT Release Summaries.
\newblock Series synthesis note (Synthesis~34), 2026. (This archive.)

\bibitem{series:downstreamnormalforms}
Anonymous.
\newblock Downstream Normal Forms for Successor Maintenance: Summary, Support, Citation, Quote, and Clause as a Stable Auditable Suffix.
\newblock Series synthesis note (Synthesis~31), 2026. (This archive.)

\bibitem{series:workedexample}
Anonymous.
\newblock Worked Example for Anonymous-DHT Receipts: Minimal Public Surface, Cross-Release Diff, and Successor Interlock.
\newblock Series synthesis note (Synthesis~17), 2026. (This archive.)

\end{thebibliography}

\end{document}
